\documentclass[10pt,a4paper]{amsart}
\usepackage[margin=19mm]{geometry}
\usepackage{amsmath,amssymb,mathtools}
\usepackage[scr=rsfs,cal=euler]{mathalfa}
\usepackage{xfrac}
\usepackage[unicode,hidelinks,bookmarks=false]{hyperref}
\newcommand{\ie}{i.e.\ }
\newcommand{\cf}{cf.\ }
\newcommand{\resp}{respectively\ }
\newcommand{\Subsection}{\subsection}
\providecommand{\nobreakdash}{\nobreak\hbox{-}}
\setlength{\emergencystretch}{2em}
\multlinegap0pt
\usepackage{amssymb}
\usepackage{xcolor}
\usepackage{mathtools}
\usepackage{tikz}
\newtheorem{theorem}{Theorem}
\def\Hc{\mathcal{H}}
\def\N{\mathbb{N}}
\newcommand{\NN}[0]{\mathbb{N}}
\newcommand{\R}{\mathcal{R}}
\title{Ramsey-like theorems and immunities}
\author{Ahmed Mimouni, Ludovic Patey}
\date{Source: arXiv:2508.15597v2 (CC BY 4.0)}
\begin{document}
\maketitle

\noindent\emph{Source and licence.} Ramsey-like theorems and immunities. Version arXiv:2508.15597v2 (CC BY 4.0), distributed by arXiv under the Creative Commons Attribution 4.0 International licence, http://creativecommons.org/licenses/by/4.0/, as recorded in the arXiv licence metadata for this exact version and shown on the abstract page https://arxiv.org/abs/2508.15597v2. Full licence notice: \texttt{licenses/arxiv-2508.15597-CC-BY-4.0.txt}.\par\smallskip

\noindent\textit{Editorial scope.} Counted unit: the article's theorem theorem (source0.tex lines 1282--1284). The article's complete original proof of it is delivered with it (source0.tex lines 1286--1316, one \texttt{proof} environment of 31 lines). No mathematics is written, repaired or re-derived: the statement and the proof are the article's own lines, in the article's own notation, and no retained line is substituted or re-flowed.  No further source-local context is needed: every cross-reference of every retained line resolves inside the two delivered spans.  Nothing was omitted from any retained span, so the package carries no omission record.  No retained line contains a citation, so the wrapper carries no bibliography and no bibliography entry.  No retained line contains a figure environment, an image-inclusion command or drawing code, so no image is delivered and the package declares no asset dependency.  Editorial formatting only, in addition: an AMS \texttt{amsart} preamble that restates the article's own declarations and macros used by the retained lines, namely the environments \texttt{theorem} on separate counters, together with the standard packages those lines need.  The retained environments print in this document as Theorem 1 in order of appearance, whereas the article prints the same items with the numbering of its own full document.  The complete native source of the article as distributed in the arXiv e-print for this exact version is bundled unchanged as \texttt{sources/183638/source0.tex}.
\begin{theorem}\label[theorem]{thm:f-2-dim-hyperimmune}
There exists a stable computable coloring $f : [\NN]^2 \to 2$ such that for each~$i < 2$, $\Hc_i(f)$ is 2-dimensional hyperimmune.
\end{theorem}

\begin{proof}
    We build the coloring $f : [\N]^2 \to 2$ by a finite injury priority argument. For
every $e \in \omega$, we want to satisfy the following requirement:
\begin{quote}
    $\R_{e,i}$ : If $\Phi_e$ is total, then there is some $n, m \in \N$ such that 
$\Phi_e(n) \subseteq A_i(f)$, $\Phi_e(n; m) \subseteq A_{1-i}(f)$ and $\Phi_e(n) \to_{1-i} \Phi_e(n; m)$.
\end{quote}

The requirements are given the usual priority ordering $\R_{0,0} < \R_{0,1} < \R_{1,0} \dots$. Initially, the
requirements are neither partially, nor fully satisfied.
\begin{itemize}
    \item A requirement $\R_{e,i}$ \emph{requires a first attention} at stage $s$ if it is not partially
satisfied and $\Phi_{e}(n)[s] \downarrow= E$ for some set $E \subseteq \{e + 1, \dots , s - 1\}$ such that
no element in E is restrained by a requirement of higher priority. If it
receives attention, then it puts a restraint on $E$, commit the elements of $E$
to be in $A_{1-i}(f)$, and is declared partially satisfied.
    \item  A requirement $\R_{e,i}$ \emph{requires a second attention} at stage $s$ if it is not fully
satisfied, $\Phi_{e}(n)[s] \downarrow= E$ and $\Phi_{e}[s](n; m) \downarrow= F$ for some sets $E < F \subseteq
\{e + 1, \dots , s - 1\}$ such that $E \to_{1-i} F$ and which are not restrained by
a requirement of higher priority. If it receives attention, then it puts a
restraint on $E \cup F$, commits the elements of $E$ to be in $A_{i}(f)$, the elements
of $F$ to be in $A_{1-i}(f)$, and is declared fully satisfied.
\end{itemize}
At stage $0$, we let $f = \emptyset$. Suppose that at stage $s$, we have defined $f(x, y)$ for every
$x < y < s$. For every $x < s$, if it is committed to be in some $A_i(f)$, set $f(x, s) = i$. Let $\R_{e,i}$ be the requirement of highest priority which
requires attention. If $\R_{e,i}$ requires a second attention, then execute the second
procedure, otherwise execute the first one. In any case, reset all the requirements
of lower priorities by setting them unsatisfied, releasing all their restraints, and go
to the next stage. This completes the construction. One easily sees by induction that each requirement $\R_{e,i}$ acts finitely often, and either $\Phi_e$ is partial, in which case $\R_{e,i}$ is vacuously satisfied, or is eventually fully satisfied. This
procedure also yields a stable coloring. 
\end{proof}

\end{document}
